\problemname{Dandelion Chaos}
\noindent

Before Gösta died of old age, scientists discovered how to make people immortal. How this was achieved
is left as an exercise for the reader. Gösta has a meadow of dandelions. After a long wait, it is finally
fully grown, and he can now create dandelion soda with the dandelions (see \textit{Dandelions}). Interestingly,
not all dandelions are exactly the same: he has discovered that there are $N$ different species of dandelions
in his meadow. Naturally, he now wants to know if certain dandelion species create tastier soda than others.
After some effort, Gösta has managed to pick exactly $2$ specimens of each species. Just as he was about to
begin analyzing them, he accidentally stumbles and mixes all the dandelions together. Disaster!

Since the dandelions are very visually similar, he will need to use his pollen analysis machine. It works by
allowing all the dandelions in it to bloom, and then it can analyze the pollen. It can then determine the
number of unique species of dandelions present. Unfortunately, it cannot dintinguish whether there is one or two
of a particular species, as the amount of pollen can vary greatly from one dandelion to another. He has numbered
all his dandelions from $1$ to $2N$.

There are two ways in which you can use the machine:
\begin{itemize}
  \item Insert a dandelion into the machine.
  \item Remove a dandelion from the machine (that is already inside of it).
\end{itemize}

After you insert or remove a dandelion from the machine, you immediately find out the number of unique species
of dandelions present in it. Gösta now wants your help to use the machine to find which pairs of dandelions
are of the same species. Since he is immortal, he is not in a hurry, but still wants to finish within a
"reasonable" time. Therefore, he can accept that you use the machine at most one million ($10^6$) times.


\section*{Implementation}
Your solution to the problem must be written in C++. It should include the file \textbf{maskroskaos.h}. 
You shall \textbf{not} create a main function. Instead, you shall implement the function\\

\begin{itemize}
  \item \texttt{void Solve(int $N$)}\\
  This function is called once for each testcase, and serves as the main function of your program.
  \begin{itemize}
    \item The parameter $N$ ($1 \leq N \leq 43000$) is the number of kinds of species that Gösta has collected.
  \end{itemize}
  Your program can use the following functions to solve the problem.
  \begin{itemize}
    \item \texttt{int Query(int $x$)}\\
    Where $x$ ($1 \leq X \leq 2N$) is the index of one the $2N$ dandelions. If $x$ is not present in the machine, it is inserted.
    If $x$ is already present in the machine, it is removed. The return value is the number of unique species of dandelions present
    in the machine after $x$ has been inserted or removed.
    If you call this function more than $10^6$ or call it with an $x$ that does not satisfy $1 \leq x \leq 2N$,
    you get \textbf{Wrong Answer}.
    \item \texttt{int Answer(int a, int b)}\\
    When you call this function, you claim that the dandelions with index $a$ and $b$ are of the same species.
    If you do any of the following, you will get \textbf{Wrong Answer}.
    \begin{itemize}
      \item If \texttt{Answer} is called with the same dandelion more than once. For example, the call \texttt{Answer(1,2)}
      followed by \texttt{Answer(1,3)} will give you \textbf{Wrong Answer} because $1$ is present in both.
      \item $1 \leq a,b \leq 2N$ is not satisfied.
      \item The dandelions with index $a$ and $b$ are not of the same species.
      \item You return from \texttt{solve} without having called \texttt{Answer} exactly $N$ times. 
    \end{itemize}
  \end{itemize}
\end{itemize}

\textbf{Note:} the functions \texttt{Query} and \texttt{Answer} should be called from the function \texttt{Solve}.
(Se example code in \texttt{maskroskaos.cpp})

\section*{Scoring}
Your solution will be tested on a set of test groups, each worth a number of points. 
Each test group contains a set of test cases. 
To get the points for a test group you need to solve all test cases in the test group.

\noindent
\begin{tabular}{| l | l | p{12cm} |}
  \hline
  \textbf{Group} & \textbf{Point value} & \textbf{Constraints} \\ \hline
  $1$    & $6$         & $N \leq 100$  \\ \hline
  $2$    & $25$        & $N \leq 15000$, for each pair $a$, $b$ where $a$ and $b$ are of the same species, it holds that
  $1 \leq a \leq N$ and $N+1 \leq b \leq 2N$. \\ \hline
  $3$    & $9$         & $N \leq 15000$ \\ \hline
  $4$    & $30$        & $N \leq 38000$ \\ \hline
  $5$    & $5$         & $N \leq 39000$ \\ \hline
  $6$    & $5$         & $N \leq 40000$ \\ \hline
  $7$    & $5$         & $N \leq 41000$ \\ \hline
  $8$    & $5$         & $N \leq 42000$ \\ \hline
  $9$    & $10$        & $N \leq 43000$ \\ \hline

\end{tabular}


\section*{Testing your solution}
To facilitate testing of your program, we provide a sample program and a sample judge. These can be downloaded from
the attachments at the bottom of the page.

To test the program, you can execute the following command:
\begin{itemize}
  \item g++ -o grader grader.cpp maskroskaos.cpp
\end{itemize}

This will create the file "grader", which you can then run to test the program. Note that this sample judge differs
from the judge used for the task. However, it is guaranteed that the real judge is \textbf{not} adaptive. This means
that the answer is determined before your program is run, and it cannot change it during runtime. Your program must
\textbf{not} read from or write to standard input/output (i.e. reading or writing to the terminal). Doing so
will likely cause confusing errors.


\section*{Input format of sample judge}
The sample judge will read input in the following format:
\begin{itemize}
  \item One line with with the integer $N$.
  \item $N$ lines that each contain the integers $a$ and $b$ ($1 \leq a,b \leq 2N$). This means that the dandelions with
  index $a$ and $b$ are of the same species. 
\end{itemize}

\section*{Output format of sample judge}
If your program solved the test case, the judge will write "Accepted. Queries: $q$", where $q$ is the number
of times your program called \texttt{Query}.

If your program incorrect, it will write "WA. $e$", where $e$ is the mistake your program made. If your program has
multiple errors, it will only write the first error it made.


\section*{Example interaction}

Consider the following input to the judge:
\begin{verbatim}
3
1 5
2 4
3 6
\end{verbatim}

A valid interaction that solves the problem can be as follows:

\begin{table}[htbp]
  \centering
  \begin{tabular}{|c|c|}
      \hline
      \textbf{Call} & \textbf{Return} \\
      \hline
      Solve(3) & (none) \\
      \hline
      Query(1) & 1 \\
      \hline
      Query(2) & 2 \\
      \hline
      Query(5) & 2 \\
      \hline
      Query(1) & 2 \\
      \hline
      Answer(1, 5) & (none) \\
      \hline
      Answer(2, 4) & (none) \\
      \hline
      Answer(3, 6) & (none) \\
      \hline
  \end{tabular}
\end{table}


